\documentclass[10pt,a4paper]{amsart}
\usepackage[margin=19mm]{geometry}
\usepackage{amsmath,amssymb,mathtools}
\usepackage[scr=rsfs,cal=euler]{mathalfa}
\usepackage{xfrac}
\usepackage[unicode,hidelinks,bookmarks=false]{hyperref}
\newcommand{\ie}{i.e.\ }
\newcommand{\cf}{cf.\ }
\newcommand{\resp}{respectively\ }
\newcommand{\Subsection}{\subsection}
\providecommand{\nobreakdash}{\nobreak\hbox{-}}
\setlength{\emergencystretch}{2em}
\multlinegap0pt
\usepackage{dsfont}
\usepackage{amssymb}
\usepackage{enumerate}
\newtheorem{theorem}{Theorem}
\newcommand{\ad}{\rm {ad}}
\renewcommand{\gg}{\mathfrak g}
\title{A note on compact homogeneous manifolds with Bismut parallel torsion}
\author{Fabio Podest\`a, Fangyang Zheng}
\date{Source: arXiv:2310.14002v1 (CC BY 4.0)}
\begin{document}
\maketitle

\noindent\emph{Source and licence.} A note on compact homogeneous manifolds with Bismut parallel torsion. Version arXiv:2310.14002v1 (CC BY 4.0), distributed by arXiv under the Creative Commons Attribution 4.0 International licence, http://creativecommons.org/licenses/by/4.0/, as recorded in the arXiv licence metadata for this exact version and shown on the abstract page https://arxiv.org/abs/2310.14002v1. Full licence notice: \texttt{licenses/arxiv-2310.14002-CC-BY-4.0.txt}.\par\smallskip

\noindent\textit{Editorial scope.} Counted unit: the article's theorem thm2.6 (source0.tex lines 346--348). The article's complete original proof of it is delivered with it (source0.tex lines 350--389, one \texttt{proof} environment of 40 lines). No mathematics is written, repaired or re-derived: the statement and the proof are the article's own lines, in the article's own notation, and no retained line is substituted or re-flowed.  Retained uncounted as the source-local context the proof needs: lines 212--214; lines 244--246.  Nothing was omitted from any retained span, so the package carries no omission record.  No retained line contains a citation, so the wrapper carries no bibliography and no bibliography entry.  No retained line contains a figure environment, an image-inclusion command or drawing code, so no image is delivered and the package declares no asset dependency.  Editorial formatting only, in addition: an AMS \texttt{amsart} preamble that restates the article's own declarations and macros used by the retained lines, namely the environments \texttt{theorem} on separate counters, together with the standard packages those lines need.  The retained environments print in this document as Theorem 1 in order of appearance, whereas the article prints the same items with the numbering of its own full document.  The complete native source of the article as distributed in the arXiv e-print for this exact version is bundled unchanged as \texttt{sources/148847/source0.tex}.
\begin{equation}\label{Bism}
\nabla^b_{e_i}e_j = 2\nabla^{\mbox{\tiny LC}}_{e_i}e_j -\nabla^c_{e_i}e_j, \qquad  \nabla^b_{\overline e_i}e_j =2\nabla^{\mbox{\tiny LC}}_{\overline e_i}e_j -\sum_k T^i_{jk}\overline{e}_k - \nabla^c_{\overline e_i}e_j.
\end{equation}

\begin{equation}
 {\mathfrak g} = {\mathfrak a} \oplus {\mathfrak g}_1 \oplus \cdots \oplus {\mathfrak g}_k \label{eq:decom}
 \end{equation}

\begin{theorem} \label{thm2.6}
Let $G={\mathbb C}^k \times G_1\times \cdots \times G_r$ where each $G_i$ is a connected, simply-connected simple complex Lie group. Then any left-invariant {\rm BTP} Hermitian metric on $G$ preserves the product structure, is K\"ahler flat on the Euclidean factor, and on each simple factor, two such metrics differ by a constant multiple of a $B$-isometry.
\end{theorem}

\begin{proof} Let $G$ be as above and $g$ a left-invariant Hermitian metric on $G$ that is {\rm BTP}. By (\ref{eq:decom}), we know that $g$ must be a product metric. Clearly the vector group factor corresponds to the kernel of the torsion $T$ which coincides with the kernel of the Bismut curvature $R^b$, hence is K\"ahler and flat. The main issue here is to determine the metric up to a constant multiple on each simple factor. So in the following let us assume that $G$ is a (connected and simply-connected) simple complex Lie group and $g$ a left-invariant {\rm BTP} Hermitian metric on $G$.

We want to understand the relationship between $g$ and the group structure of $G$.  Consider the Cartan-Killing form $B$ of $\gg$, which can be written as
\begin{equation} \label{eq:Killing}
B(X,Y) = \mbox{Tr}(\ad(X)\circ\ad(Y)) =  \sum_{r,s=1}^n T^r_{sX} T^s_{rY}.
\end{equation}
As $T$ is $\nabla^b$-parallel, we have $\nabla^bB =0$, so by using formulas \eqref{Bism} we obtain for any $1\leq i,j\leq n$ that
\begin{equation}\label{parB}
\left\{ \begin{split} \  \sum_{l=1}^n \big(  T^l_{ki}B(e_l,e_j) + T^l_{kj}B(e_i,e_l) \big) = 0,\\
\ \sum_{i=1}^l\big( \overline{T^i_{kl}}B(e_l,e_j) + \overline{T^j_{kl}}B(e_i,e_l)\big)  =0.
\end{split} \right.
\end{equation}
It is well known that for any given symmetric complex matrix $P$, there always exists a unitary matrix $Q$ such that $^t\!Q P Q $ is diagonal with non-negative real elements. So by a constant unitary change of $e$ if necessary, we may assume that $B= (B_{ij}) =\mbox{diag} \{ a_1, \ldots , a_n\}$ with $a_1\geq \cdots \geq a_n\geq 0$. Therefore \eqref{parB} are written as
\begin{equation*}
a_jT^j_{ik} + a_iT^i_{jk} =0 , \ \ \ \ a_jT^i_{jk} + a_iT^j_{ik} =0 , \ \ \ \ \forall \ 1\leq i,j,k\leq n.
\end{equation*}
From this we obtain for $1\leq i,j\leq n$
$$a_j^2\sum_r |T^j_{ir}|^2 = a_i^2 \sum_r |T^i_{jr}|^2 = a_i^2 \sum_r |T^j_{ir}|^2.$$
We then conclude that
\begin{equation} \label{eq:SPT3}
T^j_{i\ast} = T^i_{j\ast} =0  \ \ \ \mbox{whenever} \ a_i\neq a_j.
\end{equation}
Now suppose that the diagonal elements of $B$ are not all equal, say $a_1=\cdots =a_r > a_{r+1} \geq \cdots \geq a_n\geq 0$. Then for any $1\leq i\leq r$ and any $r+1\leq \alpha \leq n$, by (\ref{eq:SPT3}), we know that $T^i_{\alpha \ast } = T^{\alpha}_{i\ast}=0$. So if we let
$${\mathfrak g}_1={\mathbb C}\{ e_1, \ldots , e_r\}, \ \ \ \  {\mathfrak g}_2={\mathbb C}\{ e_{r+1}, \ldots , e_n\}, $$
then both ${\mathfrak g}_1$ and ${\mathfrak g}_2$ are ideals in the complex Lie algebra ${\mathfrak g}$, contradicting with the assumption that ${\mathfrak g}$ is simple. Therefore the diagonal matrix $B$ must be a constant multiple of the identity: $B=a_1I$.  As ${\mathfrak g}$ is a semisimple complex Lie algebra, we must have $a_1>0$, therefore
$$ B(e_i,e_j) =  a_1 \delta_{ij}. $$
Note that if  $\tilde{e}$ is another basis of ${\mathfrak g}$, with $\tilde{e}_i=\sum_{j=1}^nP_{ij}e_j$, then
$$ B(\tilde{e}_i,\tilde{e}_j) =  a_1 \sum_k P_{ik}P_{jk} , \ \ \ \mbox{while} \ \ g(\tilde{e}_i,\tilde{e}_j) = \sum_k P_{ik}\overline{P_{jk}} , $$
so we get the following relation between $g$ and $B$:
$$
B \,^t\!g^{-1} \overline{B} = a_1^2 g,
$$
which holds under any basis of ${\mathfrak g}$. Now suppose $h$ is another left-invariant Hermitian metric on $G$ that is {\em {\rm BTP}}. Then the above equation would hold for $h$ as well:
$$ B \,^t\!h^{-1} \overline{B} = a_1'^2h,$$
where $a_1'>0$ is another constant. Let us denote by $f: {\mathfrak g} \rightarrow {\mathfrak g}$ the isomorphism (as complex linear map) so that $h(X,Y)=\frac{a_1}{a_1'} g(f(X), Y)$ for any $X,Y \in {\mathfrak g}$. Then the above relations imply that
\begin{equation} \label{Bf}
 B(f(X), f(Y)) = B(X,Y) , \ \ \ \ \forall \ X, Y\in {\mathfrak g}.
\end{equation}
So a constant multiple of $h$ is related to $g$ by a $B$-isometry. This completes the proof of Theorem \ref{thm2.6}.
\end{proof}

\end{document}
